% Intended LaTeX compiler: pdflatex
\documentclass[a4paper,12pt]{article}
\usepackage[utf8]{inputenc}
\usepackage[T1]{fontenc}
\usepackage{graphicx}
\usepackage{longtable}
\usepackage{wrapfig}
\usepackage{rotating}
\usepackage[normalem]{ulem}
\usepackage{amsmath}
\usepackage{amssymb}
\usepackage{capt-of}
\usepackage{hyperref}
\usepackage{\string~/.emacs.d/common}
\author{mahmood}
\date{\today}
\title{modus ponens}
\hypersetup{
 pdfauthor={mahmood},
 pdftitle={modus ponens},
 pdfkeywords={},
 pdfsubject={},
 pdfcreator={Emacs 30.0.50 (Org mode 9.6.6)}, 
 pdflang={English}}
\begin{document}

\maketitle
\begin{note}
this document and others can be found on my website\\
\url{https://mahmoodsheikh36.github.io/}
\end{note}
\textbf{modus ponens} is a deductive argument form and \href{20230524202636-rule_of_inference.org}{rule of inference}. it can be summarized as "P implies Q. P is true. Therefore Q must also be true."\\[0pt]
the form of a modus ponens argument consists of two premises and a conclusion:\\[0pt]
\begin{enumerate}
\item if \(P\), then \(Q\)\\[0pt]
\item \(P\)\\[0pt]
\item therefore, \(Q\)
\end{enumerate}
\end{document}